\intro
Поиск элемента встречается при работе со списками товаров, книг и учебных
материалов. Один и тот же результат можно получить разными способами,
но количество действий при этом будет различаться.

Цель работы --- сравнить линейный и бинарный поиск в массиве целых чисел.
Для этого рассмотрены правила работы алгоритмов, получены оценки числа
шагов и выполнен небольшой вычислительный эксперимент на Python.

В работе сравнивается число просмотренных элементов, а не время выполнения.
Это позволяет получить воспроизводимый результат, который не зависит
от скорости компьютера. Один шаг означает одну итерацию основного цикла.
При этом внутри шага бинарного поиска может выполняться больше одного
сравнения, поэтому шаги не следует считать одинаковыми по стоимости.

Документ оформлен на основе шаблона СГУ~\cite{sgu}.
Программный код подключён из отдельного файла с помощью пакета
\texttt{minted}, который выполняет подсветку синтаксиса~\cite{minted}.
